\let\negmedspace\undefined
\let\negthickspace\undefined
\documentclass[journal]{IEEEtran}
\usepackage[a4paper, margin=10mm, onecolumn]{geometry}
%\usepackage{lmodern} % Ensure lmodern is loaded for pdflatex
\usepackage{tfrupee} % Include tfrupee package

\setlength{\headheight}{1cm} % Set the height of the header box
\setlength{\headsep}{0mm}  % Set the distance between the header box and the top of the text

\usepackage{gvv-book}
\usepackage{gvv}
\usepackage{cite}
\usepackage{amsmath,amssymb,amsfonts,amsthm}
\usepackage{algorithmic}
\usepackage{graphicx}
\usepackage{float}
\usepackage{textcomp}
\usepackage{xcolor}
\usepackage{txfonts}
\usepackage{listings}
\usepackage{enumitem}
\usepackage{mathtools}
\usepackage{gensymb}
\usepackage{comment}
\usepackage[breaklinks=true]{hyperref}
\usepackage{tkz-euclide} 
\usepackage{listings}
% \usepackage{gvv}                                        
\def\inputGnumericTable{}                                 
\usepackage[latin1]{inputenc}                                
\usepackage{color}                                            
\usepackage{array}                                            
\usepackage{longtable}                                       
\usepackage{calc}                                             
\usepackage{multirow}                                         
\usepackage{hhline}                                           
\usepackage{ifthen}                                           
\usepackage{lscape}
\usepackage{tikz}
\usetikzlibrary{patterns}

\begin{document}

\bibliographystyle{IEEEtran}
\vspace{3cm}

\title{12.602}
\author{ee25btech11063-vejith}

\maketitle
% \maketitle
% \newpage
% \bigskip
{\let\newpage\relax\maketitle}
\renewcommand{\thefigure}{\theenumi}
\renewcommand{\thetable}{\theenumi}
\setlength{\intextsep}{10pt} % Space between text and floats
\textbf{Question}:\\
If the following system has a non-trivial solution:
\begin{align*}
px + qy + rz &= 0, \\
qx + ry + pz &= 0,\\
rx + py + qz &= 0,
\end{align*}
then which one of the following options is \textbf{TRUE}? \hfill (CS 2016)
\vspace{0.5cm}

\begin{enumerate}
    \item[(a)] $p + q + r = 0$ \text{ or } $p = q = r$
    \item[(b)] $p + q + r = 0$ \text{ or } $p = q = -r$
    \item[(c)] $p + q + r = 0$ \text{ or } $p = q = r = 0$
    \item[(d)] $p + q + r = 0$ \text{ or } $p = -q = r$
\end{enumerate}
\textbf{Solution}:\\
Given equations are
\begin{align}
   \brak{p\hspace{0.5cm}q\hspace{0.5cm}r}\myvec{x\\y\\z}=0\\
   \brak{q\hspace{0.5cm}r\hspace{0.5cm}p}\myvec{x\\y\\z}=0\\
   \brak{r\hspace{0.5cm}p\hspace{0.5cm}q}\myvec{x\\y\\z}=0\\
   \implies \begin{pmatrix}
        p & q & r\\
        q & r & p\\
        r & p & q
    \end{pmatrix} \myvec{x\\y\\z} =\myvec{0\\0\\0}
\end{align}
let 
\begin{align}
    \Vec{A}=\begin{pmatrix}
        p & q & r\\
        q & r & p\\
        r & p & q
    \end{pmatrix} \text{ and } \Vec{X}=\myvec{x\\y\\z}\\
    \implies \Vec{A}\Vec{X}=\Vec{0}
\end{align}
The system of equations $\Vec{AX = 0}$ has non trivial solution if rank of $\Vec{A}<3$
\begin{align}
      \begin{pmatrix}
        p & q & r\\
        q & r & p\\
        r & p & q
    \end{pmatrix} &\xleftrightarrow{R_2 \rightarrow pR_2- qR_1} 
    \begin{pmatrix}
        p & q & r \\
        0 & pr-q^2 & p^2-qr \\
        r & p & q 
\end{pmatrix}\\
&\xleftrightarrow{R_3 \rightarrow pR_3- rR_1} \begin{pmatrix}
        p & q & r \\
        0 & pr-q^2 & p^2-qr \\
        0 & p^2-rq & pq-r^2 
\end{pmatrix}\\
&\xleftrightarrow{R_3 \rightarrow (pr-q^2)R_3 - (p^2-rq)R_2} \begin{pmatrix}
        p & q & r \\
        0 & pr-q^2 & p^2-qr \\
        0 & 0 & D 
\end{pmatrix}\\
\implies D=(pr-q^2)(pq-r^2)-(p^2-rq)(p^2-qr)=0\\
\implies (pr-q^2)(pq-r^2)=p^2qr-pr^3-pq^3+q^2r^2\\
\implies (p^2-rq)(p^2-qr)=p^4-2p^2qr+r^2q^2
\end{align}
From (11) and (12)
\begin{align}
    D=\brak{p^2qr-pr^3-pq^3+q^2r^2}-\brak{p^4-2p^2qr+r^2q^2}\\
    =-p^4+3p^2qr-pq^3-pr^3\\
    =-p\brak{p^3+q^3+r^3-3pqr}\\
    =-p\brak{p+q+r}\brak{p^2+q^2+r^2-pq-qr-pr}\\
 \implies D  =-\frac{1}{2}p\brak{p+q+r}\brak{\brak{p-q}^2 + \brak{q-r}^2 + \brak{p-r}^2}
\end{align}
$D=0$ means 
\begin{align}
    p+q+r=0 \hspace{1.5cm} \text{ or }\hspace{1.5cm} \brak{p-q}^2 + \brak{q-r}^2 + \brak{p-r}^2=0\\
    \implies  p+q+r=0 \hspace{3cm} \text{ or }\hspace{4.7cm} p=q=r
    \end{align}
 $\implies$ Given system of equations has non trivial solution if
 \begin{align}
     p+q+r=0 \hspace{1cm} \text{ or }\hspace{1cm} p=q=r
 \end{align}
\end{document}
